% ch01  §1 引言（原文 368–384 行）

\section{引言}

几何 Langlands 纲领由 Drinfeld 与 Laumon 发起，是 Drinfeld 关于整体函数域上 $\GL_2$ 的 Langlands 对应的工作的延续与推广。在几何理论中，基本研究对象从约化群 $G$ 的自守形式空间，变为代数曲线上 $G$-丛的模空间上的层范畴。

近年来，几何 Langlands 纲领对经典 Langlands 纲领以及若干相关的算术问题，都有富有成果的应用。传统做法是套用 Grothendieck 的层-函数字典，把几何理论中研究的层范畴「去范畴化」，得到算术理论中研究的函数空间。如上所述，Drinfeld 对 $\GL_2$ 的 Langlands 对应的处理用的就是这个办法（\cite{Dr}）。另一个著名的例子是 Ngô 对基本引理的证明（\cite{Ng}）。

在这篇综述中，我们要说明从几何理论通向算术理论的\emph{另一条}通道，它同样经由一个迹构造，但性质不同。这一方法的抽象版本在 \S\ref{S: cat trace} 讨论。应当指出，这一想法已经隐含在 V. Lafforgue 关于整体函数域上 Langlands 对应的工作（\cite{La}）中；Gaitsgory 的一篇非正式文章（\cite{Ga}）也有说明。

几何 Langlands 纲领的基石之一，是几何 Satake 等价：它在 $G$ 的仿射 Grassmann 簇上某类反常层的范畴，与 Langlands 对偶群 $\hat G$ 的有限维代数表示的范畴之间，建立张量等价。这是经典 Satake 同构（经由层-函数字典）的巨大推广，也可看作 Langlands 对偶群的一种概念性构造。

对等特征局部域上的约化群，几何 Satake 是 Lusztig、Ginzburg、Beilinson–Drinfeld 与 Mirković–Vilonen 等人工作的结果（\cite{Lu,Gi1,BD,MV}），现已在表示论、数学物理与（算术）代数几何中有许多应用。新近的工作 \cite{Zh1} 又对 $p$ 进群建立了这一等价，我们将在 \S\ref{S: GSat} 回顾。我们借机在 \S\ref{SS: MotSat} 讨论几何 Satake 的动机版本；它在某种意义上为本文的算术应用提供了一个玩具模型。接着在 \S\ref{S: cattrSat}，我们把上述抽象的迹构造用于 $p$ 进群的几何 Satake，由此得到对 Shimura 簇的上同调与轮的研究的一些应用，见 \S\ref{S:app to Sh}。最后两节基于与肖梁的合作工作 \cite{XZ}。

\medskip

\noindent\textbf{致谢。}本文源自作者在 Current Developments in Mathematics 2016 会议上的报告；作者感谢主办方的邀请。\S\ref{S: cattrSat} 与 \S\ref{S:app to Sh} 基于与肖梁（Liang Xiao）的合作工作，作者感谢他的合作。此外，作者感谢 Pavel Etingof、Dennis Gaitsgory 与刘一峰（Yifeng Liu）的有益讨论。
